\section{Related Work}
\label{sec:related_work}

Our work intersects three active areas of machine learning research: parameter-efficient multi-task serving, PAC-Bayesian learning theory, and trajectory optimization across deep neural networks.

\subsection{Parameter-Efficient Multi-Task Serving}
The rise of Low-Rank Adaptation (LoRA) \cite{hu21lora} and other Parameter-Efficient Fine-Tuning (PEFT) methods has motivated a shift towards multi-task ensembling on a shared base model. Initial approaches focused on static weight-space merging, such as Linear Averaging, Task Arithmetic \cite{mitchell80}, Git Re-Basin \cite{ainsworth2022git}, ZipIt! \cite{stoica2023zipit}, and TIES-Merging \cite{yadav23ties}, which combine specialized experts into a single static model. While computationally efficient, these methods are fundamentally limited because they yield a static, task-agnostic set of weights that cannot adapt dynamically to different test inputs.

To address this, dynamic model merging approaches have emerged, which route inputs dynamically to different adapters. These are broadly divided into two categories:
\begin{enumerate}
    \item \textbf{Weight-Space Routers:} Methods like the Linear Router and Parameter-Free Subspace Routing (PFSR) interpolate the weights of deep layers on a batch-by-batch basis. However, when deployed on heterogeneous streams where a single batch contains multiple distinct tasks, these methods average the weights of disparate experts, destroying the specialized parameter pathways. This phenomenon is known as \emph{Heterogeneity Collapse}.
    \item \textbf{Activation-Blending Routers:} Methods like SABLE and PAC-ZCA dynamically route inputs sample-by-sample, computing task-specific activations independently and then blending these activations using layer-specific routing coefficients. By operating in activation space, these methods are fully immune to Heterogeneity Collapse, but their layer-wise calibration has lacked a rigorous theoretical foundation.
\end{enumerate}

\subsection{PAC-Bayesian Learning Theory}
PAC-Bayesian generalization theory, pioneered by McAllester \cite{kearns89} and expanded by Catoni, provides rigorous generalization bounds for randomized (or Bayesian) classifiers. Unlike classical VC-dimension or Rademacher complexity bounds, PAC-Bayesian bounds are highly practical because they depend on the Kullback-Leibler (KL) divergence between a learned posterior $Q$ and a fixed prior $P$. 
In deep learning, PAC-Bayesian bounds have been successfully applied to analyze the generalization of stochastic neural networks and parameter-space temperature scaling (e.g., PAC-ZCA). However, existing applications typically calibrate a single global temperature parameter or treat layer-wise parameters as independent and identically distributed. Our work is the first to formulate a PAC-Bayesian bound specifically over deep layer-wise routing trajectories, leveraging a Markovian transition prior to capture depth-wise correlations.

\subsection{Trajectory Regularization and Deep Continuity}
The interpretation of deep neural networks as discretized continuous trajectories has a rich history, particularly popularized by Neural Ordinary Differential Equations (Neural ODEs) and residual network flow formulations. In these frameworks, representations or network weights flow smoothly across depth. Depth-wise smoothness has been shown to improve robustness, stabilize optimization, and prevent representation collapse. 
While prior works have applied heuristic smoothness penalties (e.g., ensembling coefficient constraints) to layer-wise merging trajectories, they lack a learning-theoretic justification. PAC-STM bridges this gap by demonstrating that a Gaussian random walk prior over network depth naturally and analytically derives a first-order ensembling smoothness penalty, providing a direct connection between physical depth continuity and rigorous PAC-Bayesian generalization guarantees.
